\section{Exact homological windows and adaptive widening}
\label{sec:window-integration}

Report 13 supplies an exact way to compute a prescribed interval of
unreduced Khovanov homology over $\F$ without constructing every chain
group. We integrate its window scanner and explicit mirror bound into the
current implementation, then add a bounded adaptive recognition stage.
This stage first queries degree zero, widens only while its budget permits,
and falls back to the complete selected backend if the answer is
inconclusive. The earlier adaptive cancellation strategy can also run
inside each window. Both mechanisms are optional.

\subsection{Why deleting chain groups can be exact}
Raw cube degree counts $1$-resolutions. Each remaining crossing contributes
homological support $\{0,1\}$, and delooping preserves homological degree.
Suppose the requested final raw interval is $[a,b]$ and $r$ crossings
remain. Retain the current complex only in
\begin{equation}\label{eq:integrated-window}
 [a-r-1,b+1].
\end{equation}
This is brutal truncation: groups outside the interval and maps incident
to them are deleted. The two extra degrees are guards.

\begin{lemma}
Truncation to \eqref{eq:integrated-window} preserves the final homology
in $[a,b]$ after the actual remaining crossing extensions and closure.
\end{lemma}
\begin{proof}
A current object of degree $p$ can contribute only to final degrees
$p+q$ with $0\leq q\leq r$. For a final degree
$j\in[a-1,b+1]$, every contributing $p=j-q$ belongs to
$[a-r-1,b+1]$. Thus all chain groups and maps between them in this
three-degree neighborhood survive unchanged. Homology in degree $j$
uses precisely its group and the incoming and outgoing differential.
Additive planar extension, closure, and the Khovanov functor preserve
these equalities.
\end{proof}

Neither guard can be dropped in general. The complex
$\F\xrightarrow{1}\F$ in degrees $a-1,a$ has zero $H^a$;
deleting the first group invents a class. The same example in degrees
$b,b+1$ proves the need for the other guard. Homology in a guard degree
can itself be artificial and is never returned as part of the answer.

After processing $i$ of $n$ crossings, the lower bound is
$a-(n-i)-1$ and the upper bound remains $b+1$.
At each step the code allocates only the delooped children in that band,
then constructs crossing saddles and transferred old maps only when both
endpoints survive. This is exactly the truncated extension, so the lemma
applies before allocation as well as after it. An object ceiling therefore
limits retained allocation, rather than an unnecessarily constructed full
extension. The frontier still advances when the retained complex is empty.

Unit cancellations may be interleaved arbitrarily: they are chain homotopy
equivalences, and the remaining additive extension and closure operations
preserve their homotopy equations. Applying the lemma after each extension
and this invariance after each cancellation proves correctness inductively.
The proof compares each intermediate state after its own actual suffix;
it does not require the truncated and full scanners to choose the same
pivots. An isolated object in degree $h$ may additionally be removed when
$[h,h+r]\cap[a,b]=\varnothing$, since its independent continued summand
has no homology in the target. Incident nonzero maps forbid this shortcut.

\subsection{Mirror choice and a parameterized complexity bound}
Let $s_0,s_1$ count circles in the all-zero and all-one complete smoothings.
Intersect the target with $[0,n]$ first; an empty intersection has zero
homology. Define
\[
 Q(n,k)=\sum_{j=0}^k2^j\binom nj,\qquad
 B_0=2^{s_0}Q(n,\min(n,b+1)).
\]
A prefix state with $j$ one-resolutions has at most $s_0+j$ closed circles:
complete it with zero-resolutions, then compare the resulting state with
the all-zero smoothing using $j$ surgeries, each changing the circle count
by at most one. Thus its delooping gives at most $2^{s_0+j}$ objects.
Summing over at most $\binom nj$ states bounds every retained allocation
by $B_0$. Cancellation and truncation remove physical objects and cannot
increase this bound; surviving labels extend injectively to cube labels.

If $w$ is the largest frontier size in the selected order, the ordinary
explicit coefficient representation has cost
\[
 2^{O(w+1)}(n+1)^{O(1)}(B_0+1)^{O(1)}.
\]
There are polynomially many object pairs and pivot updates per stage, and
at most $2^{w/2}$ dotted basis elements in a morphism space. This argument
counts objects actually present, so it requires no common boundary disk
or Catalan matching bound.

Mirroring reverses raw degree: the rank in original degree $h$ equals the
mirror rank in degree $n-h$. It preserves the crossing graph and the width
of the same crossing order, while interchanging $s_0$ and $s_1$.
The corresponding bound is
\[
 B_1=2^{s_1}Q(n,\min(n,n-a+1)).
\]
The implementation compares these exact integers and selects the smaller
bound, mapping returned degrees back to the original diagram. This is a
proved size-bound selection, not a prediction of observed running time.
The integers have $O(n)$ bits since $Q(n,k)\leq3^n$.
The recorded per-stage profile explicitly states whether its internal
degrees belong to the original or mirrored diagram.

For normalization let $\nu$ be the number of negative crossings returned
by \code{Diagram.signs()}; normalized degree is raw degree minus $\nu$.
The source braid word uses a different sign convention: the closure of
\code{[1]} on two strands has $\nu=1$ and raw homology in degree one.
Counting negative braid letters would select the wrong window.
If $s$ is the Seifert-circle count and $\mu=\min(\nu,n-\nu)$, surgery
comparison with the oriented smoothing gives
$s_0\leq s+\nu$ and $s_1\leq s+n-\nu$.
For the normalized band $[-d,d]$, the resulting bound is
\begin{equation}\label{eq:integrated-window-complexity}
 2^{O(w+s+\mu+d)}(n+1)^{O(\mu+d+1)}.
\end{equation}
It is quasi-polynomial when $w+s=O(\log^2(n+2))$ and
$\mu+d=O(\log(n+2))$. For a $t$-strand braid scanned in braid order,
$s=t$ and $w\leq2t$. These are bounds for the prescribed query.
No theorem here bounds a detecting radius for every nontrivial knot,
finds such a low-width presentation for every input, or turns an agreeing
partial query into a complete unknot recognizer.

\subsection{An introspective recognition policy}
The new recognition options are \code{window\_radius},
\code{window\_max\_objects}, and \code{window\_seconds}; the corresponding
CLI flags use hyphens. A missing radius disables the stage. An enabled
stage runs after the existing inexpensive certificates and simplification,
before the selected complete backend. Each actual factor is queried
independently and its PD is included in the evidence.

The requested radius is a maximum. The stage tries normalized bands
$[-d,d]$ at radii $0,1,2,4,\ldots$, including the requested endpoint
once and stopping when every raw degree is covered. Each wider query
restarts: discarded chain groups cannot be recovered from a narrow scan.
The last successful crossing order is reused, including when the selected
mirror changes. All queries for one factor share the same time allowance
(default $0.1$ seconds) and retained-object ceiling (default $20\,000$),
both bounded by the enclosing recognition limits. Before another query,
the policy stops if the remaining local time is less than twice the last
query's time. This is an empirical progress rule, not a competitive bound.

An exact queried profile differing from $\{0:2\}$ proves knottedness,
since every queried band includes normalized degree zero. Agreement is
inconclusive unless the query covers all raw degrees $0,\ldots,n$.
Only in that complete case does rank two certify the unknot using the
existing Khovanov detection theorem. For example, a trefoil agrees in
normalized degree zero and is detected after widening. A local time,
object, or memory limit abandons the probe and continues the full backend;
exhausting the enclosing deadline returns \code{UNKNOWN}.

This differs operationally from adaptive cancellation: that method pauses
between complete pivots and resumes the same complex when its residue
shortcut fails. Window widening must restart. Both fit the same design
principle of measuring progress, switching only at mathematically valid
boundaries, and retaining an exact fallback. Their combination does not
improve the general exponential worst-case guarantee.

The standalone \code{window} CLI and the Python functions in
\path{fastunknot/window_scan.py} and \path{fastunknot/window_bounds.py}
expose exact partial ranks independently of recognition. CLI
\code{--normalized} interprets input endpoints as normalized degrees and
adds a separate normalized profile; \code{by\_degree} always uses original
raw degrees. The API result distinguishes \code{window\_rank} from
\code{complete\_rank}. A partial rank of two alone is never a verdict.

\subsection{Validation and measured tradeoffs}
The integrated suite passes 245 tests. New checks compare 100 deterministic
random knot braids on two through five strands with full homology, using
shuffled crossing orders and arbitrary intervals. Every interval is checked
with ordinary, residue, and adaptive cancellation; the first 20 also force
the dense component kernel. Intermediate differentials are checked to
square to zero. Other regressions cover both guards, mirror degree mapping,
wrong-sign hazards, factor evidence, partial agreement, complete coverage,
local versus global limits, and budget-sensitive widening.
Separately, report 13 passes its 64 scanner tests and 10 hierarchy tests;
report 14 passes 94 tests. These later review runs do not alter the
as-delivered archive trees. Logs are retained under \path{synthesis/data/}.

\begin{center}
\begin{tabular}{@{}lrrrr@{}}
\toprule
Input & Full/window & Probe 0 & Probe 4 & A/A \\
\midrule
conway & 0.973 & 0.988 & 0.959 & 0.950 \\
kinoshita\_terasaka & 1.142 & 1.030 & 1.030 & 0.984 \\
hard\_unknot\_8 & 0.983 & 0.514 & 0.207 & 1.032 \\
stress\_braid5\_36 & 1.374 & 1.478 & 1.469 & 0.988 \\
torus\_4\_7 & 8.460 & 0.885 & 1.230 & 1.009 \\
two\_strand\_201 & 2.733 & 0.741 & 0.946 & 1.019 \\
\bottomrule
\end{tabular}
\end{center}
Ratios are median paired baseline/new times; values above one are faster.


The table uses seven shuffled paired rounds on CPython 3.13.14. A/A compares
identical full-decision arms. Preparation and imports are excluded; order
selection and mirror selection are included. The raw-window ratio compares
\emph{different output contracts}: degree zero versus all homology. The two
probe columns compare complete decision procedures with fallback, using a
two-second shared local window budget and maximum radii zero and four.
Earlier certificates are disabled in those timed arms to isolate this stage.
With default certificates enabled, Jones decides Conway and Kinoshita--Terasaka,
the short-braid engine decides the hard unknot, Alexander decides the stress
braid, and Bennequin decides both torus examples. Thus these are fallback
measurements, not claimed default-pipeline speedups.

The stress input reduces peak retained allocation from $17\,694$ to
$7\,692$ objects and its degree-zero rank is $698$, already an obstruction.
At a $10\,000$-object cap, the full baseline returns \code{UNKNOWN}; the
window stage returns \code{KNOTTED}. Conversely, the hard unknot gains no
allocation reduction and maximum-radius-four probing is about $4.8$ times
slower than the complete baseline. The large raw-window gain on the
four-strand torus example does not directly translate to a decision gain:
degree zero agrees with the unknot and requires widening or fallback.
These negative results justify keeping the policy optional.

\subsection{Review boundaries for the adjacent reports}
Report 13's finite nilpotent transfer and terminal patterned-ball routine
are separate proposals; their passing tests do not supply the missing
complete hierarchy construction. They are not substituted for a knot
recognition backend here. Report 14's interval normalization, decision
quotients, and scalar Fitting split were initially held for grading review.
Section~\ref{sec:continuation-integration} records their subsequent optional
integration, including the additional restriction of scalar basis changes
to recovered quantum-shift blocks. These operations remain separate from
the exact window scanner; an agreeing partial window may fall back to the
new complete decision backend, but its shortened model is never used to
answer an exact window query.
